% GENERATED FILE. Edit canonical lessons, then run npm run generate:books.
% canonical-source-hash: fnv1a64:9a86e7b6312263a2
% canonical-lessons: GU-C213-havaman, GU-C213-tadko, GU-C213-thandu, GU-C213-vadalchayu, GU-C213-varasvu

\chapter{Tomorrow's weather}
\label{ch:gu-tomorrow-s-weather}
\hlchaptermodality{213}

\begin{chapteropening}
\textbf{By the end of this chapter:} \emph{I can say what the weather will be like tomorrow.}

Complete the last lesson of chapter 213: I can say what the weather will be like tomorrow.
\end{chapteropening}

\section[\texorpdfstring{\gu{હવામાન} (havāmān)}{હવામાન (havāmān)}]{\gu{હવામાન} (havāmān) — weather, climate}

\label{lesson:GU-C213-havaman}



\begin{quote}
Before the new one: say the Gujarati for to believe, then the Gujarati for to change.
\end{quote}

\subsection*{You'll want to know: \gu{હવામાન}}
\textbf{\gu{હવામાન}} — \emph{havāmān} — \textquotedblleft{}weather, climate\textquotedblright{}.

\begin{quote}
\textbf{\gu{કાલે} \gu{હવામાન} \gu{સારું} \gu{હશે}.} — \emph{kāle havāmān sārũ haśe} — \textquotedblleft{}The weather will be good tomorrow.\textquotedblright{}
\end{quote}

\begin{cousinweb}
\textbf{\gu{હવામાન}} (\emph{havāmān}) is built from two older words: \emph{havā}, \textquotedblleft{}air\textquotedblright{}, a loan from Arabic \emph{hawā'}, which reached Gujarati through Persian (Lane, under the root \emph{h-w-y}), and \emph{mān}, \textquotedblleft{}measure\textquotedblright{}, from Sanskrit \emph{māna}. CDIAL has no entry for the compound.
\end{cousinweb}

\begin{grammarlens}[title={What you've built}]
One more word for everyday life at A2.
\end{grammarlens}

\subsection*{Your turn}
\begin{itemize}
\raggedright
  \item \emph{Say it:} \emph{havāmān}
  \item \emph{Say it:} \emph{havāmān}, once more
  \item \emph{Say it:} \emph{kāle havāmān sārũ haśe} — The weather will be good tomorrow.
  \item \emph{Recall:} say the Gujarati for to believe, then the Gujarati for to change, then say \emph{havāmān} again
\end{itemize}

\begin{tcolorbox}[breakable,colback=teal!4,colframe=teal!35!black,title={Before you move on}]
What does \gu{હવામાન} mean? (\textquotedblleft{}Weather, climate\textquotedblright{}.) Say it once more.
\end{tcolorbox}
\section[\texorpdfstring{\gu{તડકો} (taḍko)}{તડકો (taḍko)}]{\gu{તડકો} (taḍko) — sunshine, the sun's heat}

\label{lesson:GU-C213-tadko}



\begin{quote}
Before the new one: say the Gujarati for to change, then the Gujarati for weather.
\end{quote}

\subsection*{You'll want to know: \gu{તડકો}}
\textbf{\gu{તડકો}} — \emph{taḍko} — \textquotedblleft{}sunshine, the sun's heat\textquotedblright{}.

\begin{quote}
\textbf{\gu{આવતીકાલે} \gu{તડકો} \gu{હશે}.} — \emph{āvtīkāle taḍko haśe} — \textquotedblleft{}It will be sunny tomorrow.\textquotedblright{}
\end{quote}

A masculine noun.

\begin{cousinweb}
The origin of \textbf{\gu{તડકો}} (\emph{taḍko}) is not settled. CDIAL lists no Gujarati form like it; its nearest entry, Sanskrit \emph{tāpa}, \textquotedblleft{}heat\textquotedblright{}, gives Old Gujarati \emph{tāvaḍu}, \textquotedblleft{}heat of the sun\textquotedblright{}, but no link to \emph{taḍko} is made there (CDIAL 5767).
\end{cousinweb}

\begin{grammarlens}[title={What you've built}]
One more word for everyday life at A2.
\end{grammarlens}

\subsection*{Your turn}
\begin{itemize}
\raggedright
  \item \emph{Say it:} \emph{taḍko}
  \item \emph{Say it:} \emph{taḍko}, once more
  \item \emph{Say it:} \emph{āvtīkāle taḍko haśe} — It will be sunny tomorrow.
  \item \emph{Recall:} say the Gujarati for to change, then the Gujarati for weather, then say \emph{taḍko} again
\end{itemize}

\begin{tcolorbox}[breakable,colback=teal!4,colframe=teal!35!black,title={Before you move on}]
What does \gu{તડકો} mean? (\textquotedblleft{}Sunshine, the sun's heat\textquotedblright{}.) Say it once more.
\end{tcolorbox}
\section[\texorpdfstring{\gu{ઠંડું} (ṭhanḍũ)}{ઠંડું (ṭhanḍũ)}]{\gu{ઠંડું} (ṭhanḍũ) — cold, cool}

\label{lesson:GU-C213-thandu}



\begin{quote}
Before the new one: say the Gujarati for weather, then the Gujarati for sunshine.
\end{quote}

\subsection*{You'll want to know: \gu{ઠંડું}}
\textbf{\gu{ઠંડું}} — \emph{ṭhanḍũ} — \textquotedblleft{}cold, cool\textquotedblright{}.

\begin{quote}
\textbf{\gu{રાત્રે} \gu{પાણી} \gu{ઠંડું} \gu{હશે}.} — \emph{rātre pāṇī ṭhanḍũ haśe} — \textquotedblleft{}The water will be cold at night.\textquotedblright{}
\end{quote}

\begin{cousinweb}
\textbf{\gu{ઠંડું}} (\emph{ṭhanḍũ}) is a tadbhava, a word inherited from Sanskrit and reshaped by sound change. CDIAL lists it, glossed \textquotedblleft{}cold\textquotedblright{}, under \emph{stabdha}, \textquotedblleft{}firmly fixed; stiff; solidified (as water)\textquotedblright{}, by way of a reshaped \emph{*ṭhaṇḍha-}, so cold is what has set hard (CDIAL 13676).
\end{cousinweb}

\begin{grammarlens}[title={What you've built}]
One more word for everyday life at A2.
\end{grammarlens}

\subsection*{Your turn}
\begin{itemize}
\raggedright
  \item \emph{Say it:} \emph{ṭhanḍũ}
  \item \emph{Say it:} \emph{ṭhanḍũ}, once more
  \item \emph{Say it:} \emph{rātre pāṇī ṭhanḍũ haśe} — The water will be cold at night.
  \item \emph{Recall:} say the Gujarati for weather, then the Gujarati for sunshine, then say \emph{ṭhanḍũ} again
\end{itemize}

\begin{tcolorbox}[breakable,colback=teal!4,colframe=teal!35!black,title={Before you move on}]
What does \gu{ઠંડું} mean? (\textquotedblleft{}Cold, cool\textquotedblright{}.) Say it once more.
\end{tcolorbox}
\section[\texorpdfstring{\gu{વાદળછાયું} (vādaḷchāyũ)}{વાદળછાયું (vādaḷchāyũ)}]{\gu{વાદળછાયું} (vādaḷchāyũ) — cloudy, overcast}

\label{lesson:GU-C213-vadalchayu}



\begin{quote}
Before the new one: say the Gujarati for sunshine, then the Gujarati for cold.
\end{quote}

\subsection*{You'll want to know: \gu{વાદળછાયું}}
\textbf{\gu{વાદળછાયું}} — \emph{vādaḷchāyũ} — \textquotedblleft{}cloudy, overcast\textquotedblright{}.

\begin{quote}
\textbf{\gu{કાલે} \gu{આકાશ} \gu{વાદળછાયું} \gu{હશે}.} — \emph{kāle ākāś vādaḷchāyũ haśe} — \textquotedblleft{}The sky will be cloudy tomorrow.\textquotedblright{}
\end{quote}

\begin{cousinweb}
\textbf{\gu{વાદળછાયું}} (\emph{vādaḷchāyũ}) is built from \emph{vādaḷ}, \textquotedblleft{}cloud\textquotedblright{}, inherited from Sanskrit \emph{vārdala}, \textquotedblleft{}rainy day\textquotedblright{} (CDIAL 11567), and \emph{chāyũ}, \textquotedblleft{}covered\textquotedblright{}, from the verb \emph{chāvũ}, \textquotedblleft{}to cover\textquotedblright{}, which CDIAL lists, glossed \textquotedblleft{}to thatch\textquotedblright{}, under Sanskrit \emph{chādayati}, \textquotedblleft{}covers, veils\textquotedblright{} (CDIAL 5018).
\end{cousinweb}

\begin{grammarlens}[title={What you've built}]
One more word for everyday life at A2.
\end{grammarlens}

\subsection*{Your turn}
\begin{itemize}
\raggedright
  \item \emph{Say it:} \emph{vādaḷchāyũ}
  \item \emph{Say it:} \emph{vādaḷchāyũ}, once more
  \item \emph{Say it:} \emph{kāle ākāś vādaḷchāyũ haśe} — The sky will be cloudy tomorrow.
  \item \emph{Recall:} say the Gujarati for sunshine, then the Gujarati for cold, then say \emph{vādaḷchāyũ} again
\end{itemize}

\begin{tcolorbox}[breakable,colback=teal!4,colframe=teal!35!black,title={Before you move on}]
What does \gu{વાદળછાયું} mean? (\textquotedblleft{}Cloudy, overcast\textquotedblright{}.) Say it once more.
\end{tcolorbox}
\section[\texorpdfstring{\gu{વરસવું} (varasvũ)}{વરસવું (varasvũ)}]{\gu{વરસવું} (varasvũ) — to pour down, to rain}

\label{lesson:GU-C213-varasvu}



\begin{quote}
Before the new one: say the Gujarati for cold, then the Gujarati for cloudy.
\end{quote}

\subsection*{You'll want to know: \gu{વરસવું}}
\textbf{\gu{વરસવું}} — \emph{varasvũ} — \textquotedblleft{}to pour down, to rain\textquotedblright{}.

\begin{quote}
\textbf{\gu{આવતીકાલે} \gu{વરસાદ} \gu{વરસશે}.} — \emph{āvtīkāle varsād varasśe} — \textquotedblleft{}It will pour with rain tomorrow.\textquotedblright{}
\end{quote}

\begin{cousinweb}
\textbf{\gu{વરસવું}} (\emph{varasvũ}) is a tadbhava, a word inherited from Sanskrit and reshaped by sound change. It comes from \emph{varṣati}, \textquotedblleft{}it rains\textquotedblright{} (CDIAL 11394).
\end{cousinweb}

\begin{grammarlens}[title={What you've built}]
One more word for everyday life at A2.
\end{grammarlens}

\subsection*{Your turn}
\begin{itemize}
\raggedright
  \item \emph{Say it:} \emph{varasvũ}
  \item \emph{Say it:} \emph{varasvũ}, once more
  \item \emph{Say it:} \emph{āvtīkāle varsād varasśe} — It will pour with rain tomorrow.
  \item \emph{Recall:} say the Gujarati for cold, then the Gujarati for cloudy, then say \emph{varasvũ} again
\end{itemize}

\begin{tcolorbox}[breakable,colback=teal!4,colframe=teal!35!black,title={Before you move on}]
What does \gu{વરસવું} mean? (\textquotedblleft{}To pour down, to rain\textquotedblright{}.) Say it once more.
\end{tcolorbox}
